\label{att:asistent_otazky}

V tejto prílohe je uvedený úplný zoznam 14 testovacích otázok použitých pri
hodnotení kvality odpovedí virtuálneho asistenta PLEXO v~sekcii
\ref{sec:test-asistent}, spolu s~odpoveďami, ktoré asistent v~teste vrátil.
Otázky sú rozdelené do troch kategórií podľa metodiky uvedenej v~kapitole
\ref{sec:assistant}: faktické otázky priamo pokryté znalostnou bázou
(10 otázok, z~toho 6 v~slovenčine a~4 v~angličtine), otázky vyžadujúce
syntézu informácií z~viacerých záznamov (2 otázky) a~otázky mimo rozsahu
znalostnej bázy (2 otázky).

\subsection*{Faktické otázky -- slovenský jazyk}

\textbf{Otázka 1:} Ako sa voláš?

\textit{Odpoveď asistenta:}
Volám sa PLEXO. Som inteligentný asistent určený na podporu obsluhy
výrobných liniek riadených PLC systémami. Pomáham pri bežnej prevádzke,
diagnostike problémov a~riešení porúch.

\bigskip
\noindent\textbf{Otázka 2:} Aký je postup mäkkého reštartu linky?

\textit{Odpoveď asistenta:}
Postup mäkkého reštartu:
\begin{itemize}
      \item Skontrolujte, či sa v~linke nenachádza žiadny produkt ani osoba.
      \item Na HMI stlačte tlačidlo „STOP" alebo „Reset".
      \item Počkajte 5 sekúnd.
      \item Stlačte tlačidlo „START".
      \item Sledujte, či linka spustí normálnu prevádzku.
\end{itemize}
Ak reštart nepomôže, zaznamenajte chybové hlásenie a~kontaktujte servisného
technika.

\bigskip
\noindent\textbf{Otázka 3:} Čo znamená chybové hlásenie ,,Motor Overload" a~ako mám
reagovať?

\textit{Odpoveď asistenta:}
,,Motor Overload" znamená preťaženie motora -- motor odoberá viac prúdu, ako
je povolené. Najčastejšou príčinou je zaseknutý produkt alebo mechanická
chyba. Postupujte takto:
\begin{itemize}
      \item Linku okamžite zastavte.
      \item Skontrolujte pás a~trasu pohybu, či nie je niečo zaseknuté.
      \item Privolajte technika -- nepokúšajte sa motor reštartovať bez kontroly.
\end{itemize}

\bigskip
\noindent\textbf{Otázka 4:} Ako mám vyčistiť svetelnú bránu?

\textit{Odpoveď asistenta:}
Čistenie senzora môže vykonať obsluha:
\begin{itemize}
      \item Zastavte linku.
      \item Suchou a~čistou handričkou (alebo stlačeným vzduchom) očistite
            šošovku senzora.
      \item Skontrolujte, či je senzor pevne uchytený a~smeruje správne.
      \item Spustite linku a~skontrolujte funkciu.
\end{itemize}
Kontrolu káblov a~výmenu senzora smie vykonávať iba technik.

\bigskip
\noindent\textbf{Otázka 5:} Aké stavy môže linka LINKA-01 mať?

\textit{Odpoveď asistenta:}
LINKA-01 môže byť v~jednom z~nasledujúcich stavov:
\begin{itemize}
      \item \textbf{Čakanie} -- linka je zapnutá, dopravník stojí, čaká na produkt.
      \item \textbf{Pohyb dopredu} -- dopravník sa pohybuje od Gate~1 ku Gate~2.
      \item \textbf{Zdržanie na Gate~2} -- dopravník stojí, odpočítava sa
            3-sekundový časovač.
      \item \textbf{Pohyb dozadu} -- dopravník sa pohybuje od Gate~2 ku Gate~1.
      \item \textbf{Výstup produktu} -- dopravník sa pohybuje dozadu ešte
            1 sekundu po detekcii Gate~1.
      \item \textbf{Chyba} -- linka sa zastavila z~dôvodu poruchy.
\end{itemize}

\bigskip
\noindent\textbf{Otázka 6:} Ako dlho zostáva produkt zdržaný na Gate~2?

\textit{Odpoveď asistenta:}
Produkt zostáva pri Gate~2 zdržaný 3 sekundy. Počas tohto času dopravník
stojí a~typicky prebieha spracovanie, skenovanie alebo kontrola produktu.
Po uplynutí tejto doby sa dopravník rozbehne smerom dozadu k~Gate~1.

\subsection*{Faktické otázky -- anglický jazyk}

\textbf{Otázka 7:} What does HMI stand for and what is it used for?

\textit{Odpoveď asistenta:}
HMI stands for Human Machine Interface. It is a~touch panel or display
through which the operator controls the production line. The HMI is one of
the basic components of a~PLC system and allows the operator to start and
stop the line, monitor its current state and view error messages.

\bigskip
\noindent\textbf{Otázka 8:} What is the purpose of the E-STOP button?

\textit{Odpoveď asistenta:}
The E-STOP (Emergency Stop) is a~red button used to immediately stop the
line in case of danger. It must be pressed whenever there is a~risk to
people or equipment. After the E-STOP has been pressed, the line must not
be restarted before the cause of the stop has been checked. The button is
released by turning it according to its type and confirming on the HMI.

\bigskip
\noindent\textbf{Otázka 9:} How does the LINKA-01 production cycle work?

\textit{Odpoveď asistenta:}
The LINKA-01 cycle consists of the following steps:
\begin{enumerate}
      \item A~product enters the belt and breaks the light beam of Gate~1.
            The PLC starts the conveyor moving forward.
      \item The product reaches Gate~2 and breaks its beam. The conveyor
            stops and a~3-second timer starts.
      \item After 3 seconds the conveyor starts moving in reverse, back
            towards Gate~1.
      \item When the product breaks Gate~1 again, the conveyor continues in
            reverse for 1 more second so the product leaves the belt.
      \item The conveyor stops and the line returns to the idle state,
            waiting for the next product.
\end{enumerate}

\bigskip
\noindent\textbf{Otázka 10:} What should I do if I see a~,,Communication Error" message?

\textit{Odpoveď asistenta:}
,,Communication Error" means that the PLC has lost connection with one of
the connected devices (sensor, motor or another module). Try a~soft restart
of the line first (STOP, wait 5 seconds, START). If the error persists,
record the message and call a~service technician. Do not attempt to fix the
issue with repeated restarts without consulting a~technician.

\subsection*{Otázky vyžadujúce syntézu}

\textbf{Otázka 11:} Produkt sa rozbehol na Gate~1, dopravník ide dopredu,
ale produkt nedorazí na Gate~2 a~celá linka sa zasekne. Čo mám robiť?

\textit{Odpoveď asistenta:}
Tento problém má niekoľko možných príčin:
\begin{itemize}
      \item Produkt mohol vypadnúť z~pásu počas presunu -- zastavte linku,
            odstráňte produkt z~okolia a~reštartujte.
      \item Gate~2 môže byť znečistená a~nedeteguje produkt, takže linka
            čaká donekonečna -- zastavte linku, suchou handričkou vyčistite
            šošovku Gate~2 a~reštartujte.
      \item Produkt môže byť príliš malý alebo ľahký na zachytenie senzorom
            -- konzultujte nastavenie s~technológom.
      \item Pás sa môže šmýkať a~produkt neprenášať -- v~tomto prípade
            privolajte technika.
\end{itemize}
Ak sa po vyčistení senzora a~mäkkom reštarte (STOP, počkať 5 sekúnd, START)
problém zopakuje, kontaktujte servisného technika.

\bigskip
\noindent\textbf{Otázka 12:} Aké kroky musím urobiť na začiatku smeny pred spustením
linky LINKA-01 a~aký je postup pri jej zastavení na konci smeny?

\textit{Odpoveď asistenta:}
\textbf{Spustenie na začiatku smeny:}
\begin{enumerate}
      \item Skontrolujte, či je pás čistý a~nie sú na ňom žiadne produkty
            ani cudzie predmety.
      \item Vizuálne skontrolujte Gate~1 a~Gate~2 -- šošovky musia byť čisté
            a~nepoškodené.
      \item Skontrolujte, či nie je stlačený E-STOP. Ak áno, uvoľnite ho otočením.
      \item Zapnite hlavný vypínač linky a~počkajte na spustenie HMI.
      \item Skontrolujte, či nie sú aktívne žiadne chybové hlásenia.
      \item Prepnite do automatického režimu a~stlačte ,,START".
\end{enumerate}
\textbf{Zastavenie na konci smeny:}
\begin{enumerate}
      \item Počkajte, kým aktuálny cyklus dokončí spracovanie a~produkt
            opustí pás.
      \item Na HMI stlačte ,,STOP".
      \item Pre dlhšiu odstávku vypnite hlavný vypínač.
\end{enumerate}

\subsection*{Otázky mimo rozsahu znalostnej bázy}

\textbf{Otázka 13:} Aká je aktuálna cena PLC Siemens S7-1200 v~Európe?

\textit{Odpoveď asistenta} (s~použitím nástroja \texttt{web\_search}):
Túto informáciu nemám priamo v~znalostnej báze, preto som ju vyhľadal na
internete. Cena PLC Siemens S7-1200 sa pohybuje rádovo v~stovkách až nízkych
tisíckach eur podľa konkrétnej konfigurácie CPU a~rozširujúcich modulov.
Pre presnú aktuálnu cenu odporúčam kontaktovať autorizovaného distribútora
Siemens. \textit{(Zdroj: vyhľadávanie cez web\_search)}

\bigskip
\noindent\textbf{Otázka 14:} What is the Modbus TCP protocol and what port does it use?

\textit{Odpoveď asistenta} (s~použitím nástroja \texttt{web\_search}):
This information is not in my knowledge base, so I~looked it up online.
Modbus TCP is an industrial communication protocol that wraps the original
Modbus messaging protocol over standard Ethernet using TCP/IP. It uses
TCP port 502. Modbus TCP follows a~client-server model and is widely used
for connecting PLCs, HMIs and SCADA systems in industrial automation.
\textit{(Source: web\_search lookup)}
